% Source: https://artofproblemsolving.com/community/c3h3537004_interesting_way_to_derive_the_quadratic_formula
% Stable identity: aops-topic-3537004-q01
If the roots are $r,s$, Vieta gives $r+s=-\dfrac{b}{a}$ and $rs=\dfrac{c}{a}$. Thus $(r-s)^2=(r+s)^2-4rs=\dfrac{b^2-4ac}{a^2}$. For $r\ge s$, one has $r-s=\dfrac{\sqrt{b^2-4ac}}{|a|}$. Adding and subtracting this from the root sum gives $r,s=-\dfrac{b}{2a}\pm\dfrac{\sqrt{b^2-4ac}}{2|a|}$. As an unordered pair, these are exactly $\boxed{\dfrac{-b\pm\sqrt{b^2-4ac}}{2a}}$. The absolute value matters only in identifying which sign labels the larger root.
